% --- Chapter appendices, numbered <chapter>.A, .B, ... -----------------------
\renewcommand{\thesection}{\thechapter.\Alph{section}}
\setcounter{section}{0}

\section{Checks and limiting cases}
\label{bdc:app:limits}

The main derivations assume positive birth and catastrophe rates and begin from
one founder. This appendix checks the endpoints of the closed forms and then
approaches the boundaries at which the two-root description degenerates.

\subsection{Initial values}
\label{bdc:app:initial}

At $t=0$, the process contains one unit and no event has occurred. Thus
\[
I(0)=1,\qquad D(0)=0,\qquad \Ifix(0)=1,\qquad
J(0)=K(0)=1,\qquad V(0)=0.
\]
Setting $w=1$ recovers each value from the corresponding closed form. The
initial slopes also have a direct interpretation. Equation~\cref{bdc:eq:IJ}
gives
\[
I'(0)=-\delta J(0)=-\delta,
\]
the rate at which the founder triggers rupture. Similarly,
\[
\Ifix'(0)=-(\lambda+\mu+\delta)+\lambda=-(\mu+\delta).
\]
Birth is absent from this loss rate because division leaves the cell productive;
only death or rupture can remove it at the first instant.

\subsection{Late-time values}
\label{bdc:app:late}

As $t\to\infty$, $w\to\infty$, and
\[
I\to b,\qquad D\to b,\qquad \Ifix\to0,\qquad
J\to0,\qquad K\to0,\qquad
V\to V_\infty=\frac{aB}{A}.
\]
The equality of the first two limits is not accidental. At late time, a path
which has avoided rupture must have died internally; the productive probability
is the vanishing difference between those two events.

\subsection{No death}
\label{bdc:app:nodeath}

Set $\mu=0$. From $ab=\mu/\lambda$ we have $b=0$, while
$a=(\lambda+\delta)/\lambda$. Internal extinction is then impossible:
$D\equiv0$ and $\Ifix=I$. Every path ruptures eventually, and
\[
V_\infty=1+\frac{\lambda}{\delta}
=\ex{\Kb\mid\text{burst}}.
\]
This is the simplest setting in which to check the geometric burst-size moments,
as done in \cref{bdc:sec:conditioning}.

\subsection{No catastrophe}
\label{bdc:app:nocatastrophe}

Setting $\delta=0$ removes rupture and returns the classical linear
birth--death process. Since the main formulae contain $\delta$ in their
denominators, we take the limit $\delta\downarrow0$ in the supercritical regime
$\lambda>\mu$. Then
\[
a\to1,\qquad b\to\frac{\mu}{\lambda},\qquad
A=a-1\to0,\qquad I(t)\to1.
\]
The internal-extinction probability approaches
\[
D(t)\longrightarrow
\frac{\mu}{\lambda}
\frac{w-1}{w-\mu/\lambda},
\qquad
w=\ee^{(\lambda-\mu)t},
\]
which is the standard finite-time extinction probability of the linear
birth--death process from one founder. At
$(\lambda,\mu,t)=(1,0.3,3)$, the limiting formula gives $0.27330340$.
Meanwhile $V_\infty\to\infty$: for positive but increasingly small
$\delta$, a supercritical population can grow for an increasingly long time
before release.

\subsection{Critical births and deaths}
\label{bdc:app:critical}

When $\lambda=\mu$, Vieta's relation gives $ab=1$, and hence
\[
b=\frac{1}{a}.
\]
At criticality, the probability of internal extinction before rupture equals
the ratio of the geometric burst-size law. For
$(\lambda,\mu,\delta)=(1,1,0.05)$, the two roots are $a=1.25$ and $b=0.8$.
Away from criticality the equality generally disappears.

\subsection{No births}
\label{bdc:app:nobirths}

Finally, let $\lambda\downarrow0$. The upper root escapes to infinity and the
quadratic Riccati equation becomes linear:
\[
\ddt{I}=\mu-(\mu+\delta)I,
\qquad I(0)=1.
\]
The remaining finite root tends to
\[
b\longrightarrow\frac{\mu}{\mu+\delta}.
\]
This is exactly the chance that a non-replicating founder dies before it
triggers rupture: the two exponential clocks divide the outcome in proportion
to their rates. The two-root machinery has degenerated, but its surviving root
retains the correct probabilistic meaning.

\subsection{The second moment of the release}
\label{bdc:app:Wsq}

At rupture from state $n$, $W_t^2$ increases from zero to $n^2$, and this
happens at rate $\delta n$. Therefore
\begin{equation}
\label{bdc:eq:ddWsq}
\ddt{}\ex{W_t^2}=\delta\ex{X_t^3}.
\end{equation}
The third moment has appeared, but the equation for $K$ contains the same term.
A unit change in $X_t$ does not produce a unit change in $X_t^2$.

\begin{table}[!htbp]
\centering
\caption{Changes in $X_t^2$ under the three event types.}
\label{bdc:tab:xsqchanges}
\begin{tabular}{lcl}
\toprule
Rate & Change in $X_t$ & Change in $X_t^2$\\
\midrule
$\lambda X_t$ & $+1$ & $(X_t+1)^2-X_t^2=2X_t+1$\\
$\mu X_t$ & $-1$ & $(X_t-1)^2-X_t^2=-2X_t+1$\\
$\delta X_t$ & $-X_t$ & $(X_t-X_t)^2-X_t^2=-X_t^2$\\
\bottomrule
\end{tabular}
\end{table}

The entries in \cref{bdc:tab:xsqchanges} give
\begin{align*}
\ex{X_{t+\Delta t}^2-X_t^2}
&=\ex{\lambda X_t\Delta t\,(2X_t+1)
+\mu X_t\Delta t\,(-2X_t+1)
+\delta X_t\Delta t\,(-X_t^2)}\\
&=\Bigl[2(\lambda-\mu)\ex{X_t^2}
+(\lambda+\mu)\ex{X_t}
-\delta\ex{X_t^3}\Bigr]\Delta t,
\end{align*}
and hence
\begin{equation}
\label{bdc:eq:ddK}
\ddt{K}
=2(\lambda-\mu)K+(\lambda+\mu)J-\delta\ex{X_t^3}.
\end{equation}
Using \cref{bdc:eq:ddWsq} to remove the third moment,
\begin{equation}
\label{bdc:eq:dWsq}
\ddt{}\ex{W_t^2}
=2(\lambda-\mu)K+(\lambda+\mu)J-\ddt{K}.
\end{equation}
Integration from $0$ to $t$, together with $V'=\delta K$ and
$I'=-\delta J$, gives
\begin{align*}
\ex{W_t^2}
&=2(\lambda-\mu)\int_0^tK\,\dd s
+(\lambda+\mu)\int_0^tJ\,\dd s-K(t)+K(0)\\
&=\frac{2(\lambda-\mu)}{\delta}V(t)
+\frac{\lambda+\mu}{\delta}[1-I(t)]-K(t)+K(0),
\end{align*}
where $K(0)=1$. This is \cref{bdc:eq:Wsq}.
